\section{Markov specifications on the integers}

Throughout this section the parameter set is $S = \mathbb{Z}$ and the state space $E$ is a
finite non-empty set with $\mathcal{E}$ its power set. Georgii takes the a priori measure
$\lambda$ on $E$ to be counting measure; the Lean development uses the uniform probability
measure on $E$ instead. By Remark (1.28)(3) the two choices give the same modifications and hence
the same Gibbs measures, and the normalising constants differ by the explicit factor
$|E|^{|\Lambda|}$ which is visible in the Lean partition-function lemmas.

For $\Lambda \subseteq \mathbb{Z}$ Georgii writes
\[
    \partial\Lambda = \set{i \in \mathbb{Z}\setminus\Lambda : |i - j| = 1 \text{ for some }
    j \in \Lambda} \qquad (3.4)
\]
for the boundary of $\Lambda$. There is no Lean definition of $\partial\Lambda$ for a general
finite $\Lambda$: every finite-volume statement below is proved for an interval
$\Lambda = [a,b]$, whose boundary is $\set{a-1, b+1}$, and the two boundary sites appear
explicitly in the formulas.

One further difference in presentation. Georgii \emph{defines} the specification attached to a
stochastic matrix $P$ by the conditioning equation (3.6) and only afterwards identifies it as
Gibbsian for a nearest-neighbour potential (Corollary (3.9)). The Lean development runs the other
way: $\gamma_P$ is \emph{defined} as the Gibbsian specification of the potential
$\Phi_{\set{i,i+1}} = -\log P(\sigma_i, \sigma_{i+1})$, and (3.6), (3.8)(1) and (3.11) are then
theorems about it. Nothing is lost, and the properness and consistency of $\gamma_P$ come for
free from the Gibbsian machinery instead of Georgii's ad hoc argument in step 1.

\begin{definition}[Positive homogeneous Markov specification, Georgii (3.1)]
    \label{def:mkv-pos-hom-markov}
    \uses{def:spec}
    \lean{MeasureTheory.GibbsMeasure.Markov.IsPositiveHomogeneousMarkovWith, MeasureTheory.GibbsMeasure.Markov.IsPositiveHomogeneousMarkov}
    \leanok

    A specification $\gamma$ with parameter set $\mathbb{Z}$ and state space $E$ is a
    \textbf{positive homogeneous Markov specification} if there is a function $g > 0$ on $E^3$
    with
    \[
        \gamma_{\set{i}}(\sigma_i = y \mid \omega) = g(\omega_{i-1}, y, \omega_{i+1})
    \]
    for all $i \in \mathbb{Z}$, $y \in E$ and $\omega \in \Cfg$. In Lean the notion is split in
    two: a structure carrying the strict positivity of $g$ together with the displayed identity
    for the singleton kernels, and its existential closure over $g$.
\end{definition}

\begin{remark}[Comment (3.2): the determining function]
    \label{rmk:mkv-comment-32}
    \uses{def:mkv-pos-hom-markov, def:spec}
    \lean{Specification.eq_of_forall_singleton_eq, MeasureTheory.GibbsMeasure.Markov.eq_markovSpecification_of_forall_singleton}
    \leanok

    Definition \ref{def:mkv-pos-hom-markov} constrains only the singleton part of $\gamma$, but
    by Theorem (1.33) a specification is determined by its singleton kernels, so a positive
    homogeneous Markov specification is determined by $g$, which is therefore called its
    \textbf{determining function}. The library has the general statement --- two specifications
    with equal singleton kernels agree, provided the singleton kernels are densities against a
    disjointly consistent independent family --- and the instance of it for the specifications of
    this section; both are named above.
\end{remark}

\begin{definition}[The homogeneous nearest-neighbour potential of a matrix, Georgii (3.9)]
    \label{def:mkv-potential}
    \lean{MeasureTheory.GibbsMeasure.Markov.markovPotential, MeasureTheory.GibbsMeasure.Markov.isPotential_markovPotential, MeasureTheory.GibbsMeasure.Markov.isFiniteRange_markovPotential, MeasureTheory.GibbsMeasure.Markov.isAbsolutelySummable_markovPotential}
    \leanok

    For a matrix $P$ on $E$ put
    \[
        \Phi^P_A = \begin{cases}
            -\log P(\sigma_i, \sigma_{i+1}) & \text{if } A = \set{i, i+1},\\
            0 & \text{otherwise.}
        \end{cases}
    \]
    This is a homogeneous nearest-neighbour potential in Georgii's sense (with
    $\varphi_1 = 0$ and $\varphi_2 = -\log P$). It is a potential (each $\Phi^P_A$ is
    $\mathcal{F}_A$-measurable), it has finite range (the interactions containing $i$ lie inside
    $[i-1, i+1]$), and it is absolutely summable: $|\Phi^P_A| \le \sum_{x,y}|\log P(x,y)|$
    uniformly, and only three sets $A \ni i$ contribute. Hence $\Phi^P \in \Bspace$.
\end{definition}

\begin{definition}[The Markov specification of a matrix, Georgii (3.6)]
    \label{def:mkv-specification}
    \uses{def:mkv-potential, def:spec, def:isssd, def:modif}
    \lean{MeasureTheory.GibbsMeasure.Markov.markovSpecification}
    \leanok

    For a positive stochastic matrix $P$, $\gamma_P$ is the Gibbsian specification of the
    potential $\Phi^P$ of Definition \ref{def:mkv-potential} at inverse temperature $1$, with the
    uniform probability measure on $E$ as a priori measure: $\gamma_P$ is the modification of
    $\ISSSD(\lambda)$ by the normalised Boltzmann factor
    $e^{-H_\Lambda^{\Phi^P}} / Z_\Lambda$.
\end{definition}

\begin{lemma}[Perron--Frobenius, Georgii Appendix 3.A and Comment (3.8)(2)]
    \label{lem:mkv-perron}
    \lean{Matrix.exists_unique_perron, Matrix.perronRoot, Matrix.perronVector, Matrix.mulVec_perronVector, Matrix.exists_eq_smul_perronVector}
    \leanok

    A matrix $Q$ on $E$ with all entries strictly positive has a strictly positive eigenvalue
    $q$ admitting a strictly positive right eigenvector $r$; $q$ is the only eigenvalue with a
    strictly positive eigenvector, and every eigenvector for $q$ is a multiple of $r$.

    Georgii's Comment (3.8)(2) also asserts that all other eigenvalues of $Q$ have modulus
    strictly less than $q$. That spectral dominance is \emph{not} formalized. It is not needed
    here: everything below uses only the uniqueness of $q$ among eigenvalues with a positive
    eigenvector, which pins down $P = (Q(x,y)r(y)/qr(x))_{x,y}$ just as well.
\end{lemma}
\begin{proof}
    \leanok

    Collatz--Wielandt: maximise $y \mapsto \min_i (Qy)_i / y_i$ over the compact image of the
    standard simplex under the normalised map $y \mapsto Qy/\sum_i (Qy)_i$. A maximiser is a
    positive eigenvector with positive eigenvalue. If $Qv = rv$ and $Qw = sw$ with $v, w > 0$,
    comparing $v$ and $w$ at a site where $v/w$ is extremal forces $r = s$ and $v \in \mathbb{R}w$.
\end{proof}

\begin{lemma}[Doeblin's ergodic theorem, Georgii Appendix 3.A]
    \label{lem:mkv-doeblin}
    \lean{Matrix.exists_stationary, Matrix.exists_unique_stationary, Matrix.tendsto_pow_apply, Matrix.abs_pow_apply_sub_le, Matrix.pos_of_vecMul_eq_self}
    \leanok

    Let $P$ be a stochastic matrix on $E$ with strictly positive entries. Then there is a unique
    probability vector $\alpha_P$ with $\alpha_P P = \alpha_P$, it satisfies
    $\alpha_P \in \,]0,1[^E$, and $\lim_{k\to\infty} P^k(x, \cdot) = \alpha_P$ uniformly in
    $x \in E$: with $\varepsilon = \min_{x,y} P(x,y)$ and
    $\rho = 1 - |E|\varepsilon \in [0,1[$ one has $|P^k(x,y) - \alpha_P(y)| \le 2\rho^k$.
\end{lemma}
\begin{proof}
    \leanok

    Left multiplication by $P$ contracts the $\ell^1$-norm of a sum-zero vector by the Dobrushin
    factor $\rho = 1 - |E|\varepsilon < 1$. Hence $(\mu P^k)_k$ is Cauchy for any probability
    vector $\mu$; its limit is stationary, the contraction gives uniqueness and the geometric
    rate, and taking $\mu$ a point mass at $x$ gives the entrywise statement. Positivity of
    $\alpha_P$ follows from $\alpha_P = \alpha_P P$ and $P > 0$.
\end{proof}

\begin{definition}[The stationary Markov chain $\mu_P$, Georgii (3.3)]
    \label{def:mkv-chain}
    \uses{lem:mkv-doeblin, def:product-probability-measure}
    \lean{MeasureTheory.GibbsMeasure.Markov.stationaryDist, MeasureTheory.GibbsMeasure.Markov.stationaryChain, MeasureTheory.GibbsMeasure.Markov.isProbabilityMeasure_stationaryChain}
    \leanok

    For a stochastic matrix $P$ with non-vanishing entries, $\mu_P \in \Prob(\Cfg, \mathcal{F})$
    is the distribution of the stationary Markov chain with transition matrix $P$: the unique
    probability measure with
    \[
        \mu_P(\sigma_i = x_0, \dots, \sigma_{i+n} = x_n)
          = \alpha_P(x_0) P(x_0,x_1)\cdots P(x_{n-1},x_n) \qquad (3.3)
    \]
    where $\alpha_P$ is the stationary vector of Lemma \ref{lem:mkv-doeblin}. In Lean it is built
    from these finite-dimensional distributions: the interval weights are shown to form a
    projective family (the right endpoint is dropped using $\sum_y P(x,y) = 1$, the left endpoint
    using $\alpha_P P = \alpha_P$), and $\mu_P$ is their projective limit by the Kolmogorov
    extension theorem for standard Borel spaces.
\end{definition}

\begin{theorem}[Cylinder probabilities of $\mu_P$, Georgii (3.3)]
    \label{thm:mkv-chain-cylinder}
    \uses{def:mkv-chain, def:cylinder-event}
    \lean{MeasureTheory.GibbsMeasure.Markov.markovChain_cylinder}
    \leanok

    For $a \le b$ and any $\sigma \in \Cfg$,
    \[
        \mu_P\bigl(\set{\tau : \tau_k = \sigma_k \text{ for all } k \in [a,b]}\bigr)
          = \alpha_P(\sigma_a) \prod_{k=a}^{b-1} P(\sigma_k, \sigma_{k+1}).
    \]
\end{theorem}
\begin{proof}
    \leanok

    A one-point cylinder over $[a,b]$ is the preimage of a singleton under the restriction map,
    so its $\mu_P$-measure is the corresponding finite-dimensional weight by the projective limit
    property of Definition \ref{def:mkv-chain}.
\end{proof}

\begin{theorem}[Georgii, Comment (3.8)(1), one block]
    \label{thm:mkv-38-1-single}
    \uses{def:mkv-specification, def:cylinder-event}
    \lean{MeasureTheory.GibbsMeasure.Markov.markovSpecification_Icc_apply_cyl}
    \leanok

    Let $P$ be a positive stochastic matrix, $\Lambda = [a,b]$ with $b = a + n$, and let
    $\eta$ agree with $\omega$ off $\Lambda$. Then
    \[
        \gamma_{P,\Lambda}(\sigma_\Lambda = \eta_\Lambda \mid \omega)
          = \frac{P(\omega_{a-1}, \eta_a) P(\eta_a, \eta_{a+1}) \cdots P(\eta_b, \omega_{b+1})}
                 {P^{n+2}(\omega_{a-1}, \omega_{b+1})}.
    \]
    This is Georgii's Comment (3.8)(1) for a single block $\Lambda = \set{i+1, \dots, i+n_k}$.
\end{theorem}
\begin{proof}
    \leanok

    The Hamiltonian of $\Phi^P$ on $[a,b]$ is $\sum_{j=a-1}^{b} -\log P(\sigma_j,\sigma_{j+1})$,
    so the Boltzmann factor is the path weight $\prod_{j=a-1}^{b} P(\sigma_j, \sigma_{j+1})$.
    Integrating that path weight against the independent kernel on $[a,b]$ contracts the path
    into a matrix power by induction on $n$ (summing out the top site turns $g$ into $Pg$), so the
    partition function is $|E|^{-(n+1)} P^{n+2}(\omega_{a-1}, \omega_{b+1})$. The numerator is the
    single surviving term of the cylinder integral, carrying the same volume factor
    $|E|^{-(n+1)}$, which cancels.
\end{proof}

\begin{theorem}[Georgii, Comment (3.8)(1), one block inside an interval]
    \label{thm:mkv-38-1}
    \uses{def:mkv-specification, thm:mkv-38-1-single, def:cylinder-event}
    \lean{MeasureTheory.GibbsMeasure.Markov.markovSpecification_Icc_apply_cyl_of_subset}
    \leanok

    Let $\Lambda = [l, m]$ with $m = l + d$ sit inside $\Delta = [a, b]$ with $l - 1 = a + p$ and
    $b = m + 1 + q$. Then for all $\zeta, \omega \in \Cfg$,
    \[
        \gamma_{P,\Delta}(\sigma_\Lambda = \zeta_\Lambda \mid \omega)
          = \frac{P^{p+2}(\omega_{a-1}, \zeta_l)\, P(\zeta_l, \zeta_{l+1}) \cdots
                  P(\zeta_{m-1}, \zeta_m)\, P^{q+2}(\zeta_m, \omega_{b+1})}
                 {P^{d+p+q+4}(\omega_{a-1}, \omega_{b+1})}.
    \]
    Georgii's Comment (3.8)(1) computes $\mu_P(\sigma_\Lambda = \zeta \mid \sigma_{\partial\Lambda}
    = \omega_{\partial\Lambda})$ for a $\Lambda$ that is a finite union of blocks, with
    conditioning volume $\Lambda$ itself. The Lean statement is at different hypotheses: one block
    $\Lambda$ inside one interval $\Delta$, with $\Delta$ arbitrary. The multi-block case is not
    formalized; the case $\Delta = \Lambda$ is Theorem \ref{thm:mkv-38-1-single}.
\end{theorem}
\begin{proof}
    \leanok

    Split $\Delta$ into the two outer blocks $[a, l-1]$, $[m+1, b]$ and the middle block
    $\Lambda$. Integrating the path weight over the outer blocks contracts each of them into a
    matrix power, by the same induction as in Theorem \ref{thm:mkv-38-1-single}, leaving the
    middle path weight untouched; the cylinder indicator then pins the middle block to $\zeta$.
    Dividing by the partition function of $\Delta$ cancels the volume factors.
\end{proof}

\begin{theorem}[The determining function of $\gamma_P$, Georgii (3.11)]
    \label{thm:mkv-311}
    \uses{thm:mkv-38-1-single, def:mkv-pos-hom-markov, def:mkv-specification}
    \lean{MeasureTheory.GibbsMeasure.Markov.markovDeterminingFun, MeasureTheory.GibbsMeasure.Markov.markovSpecification_singleton_apply, MeasureTheory.GibbsMeasure.Markov.isPositiveHomogeneousMarkovWith_markovSpecification, MeasureTheory.GibbsMeasure.Markov.isPositiveHomogeneousMarkov_markovSpecification}
    \leanok

    For a positive stochastic matrix $P$,
    \[
        \gamma_{P,\set{i}}(\sigma_i = y \mid \omega)
          = g_P(\omega_{i-1}, y, \omega_{i+1}), \qquad
        g_P(x,y,z) = \frac{P(x,y)P(y,z)}{P^2(x,z)} \qquad (3.11)
    \]
    so $\gamma_P$ is a positive homogeneous Markov specification with determining function $g_P$.
\end{theorem}
\begin{proof}
    \leanok

    Theorem \ref{thm:mkv-38-1-single} with $\Lambda = [i,i]$, i.e.\ $n = 0$: the numerator is
    $P(\omega_{i-1}, y)P(y, \omega_{i+1})$ and the partition function is
    $P^2(\omega_{i-1},\omega_{i+1})$. Positivity of $g_P$ is positivity of $P$.
\end{proof}

\begin{definition}[The matrix attached to a determining function, Georgii (3.7)]
    \label{def:mkv-matrix-of-det-fun}
    \uses{lem:mkv-perron, def:mkv-pos-hom-markov}
    \lean{MeasureTheory.GibbsMeasure.Markov.detQ, MeasureTheory.GibbsMeasure.Markov.matrixOfDetFun}
    \leanok

    Let $g : E^3 \to \,]0,\infty[$ and fix $a \in E$. Put
    $Q(x,y) = g(a,x,y)/g(a,a,y)$ and let $q > 0$, $r \in \,]0,\infty[^E$ be the Perron root and a
    Perron eigenvector of $Q$ (Lemma \ref{lem:mkv-perron}). Then
    \[
        P(x,y) = \frac{Q(x,y) r(y)}{q\, r(x)}. \qquad (3.7)
    \]
\end{definition}

\begin{proposition}[Comment (3.8)(2): (3.7) is well defined]
    \label{prop:mkv-38-2}
    \uses{def:mkv-matrix-of-det-fun, lem:mkv-perron}
    \lean{MeasureTheory.GibbsMeasure.Markov.matrixOfDetFun_pos, MeasureTheory.GibbsMeasure.Markov.matrixOfDetFun_mem_rowStochastic, MeasureTheory.GibbsMeasure.Markov.matrixOfDetFun_eq_of_eigen}
    \leanok

    Equation (3.7) is a well-defined map from the positive functions $g$ on $E^3$ to the positive
    stochastic matrices on $E$: the matrix of Definition \ref{def:mkv-matrix-of-det-fun} has
    strictly positive entries, its rows sum to $1$, and it does not depend on the choice of
    Perron eigenvector $r$ (any positive eigenpair $(q', r')$ of $Q$ gives the same $P$).
\end{proposition}
\begin{proof}
    \leanok

    Positivity is clear from $Q > 0$, $r > 0$, $q > 0$. For stochasticity, sum (3.7) over $y$ and
    use $Qr = qr$. Independence of the choice: by Lemma \ref{lem:mkv-perron} a positive
    eigenvector is unique up to a positive scalar and its eigenvalue is $q$, and the scalar
    cancels between $r(y)$ and $r(x)$ in (3.7).
\end{proof}

\begin{theorem}[Georgii, step 2 of the proof of (3.5)]
    \label{thm:mkv-step2}
    \uses{thm:mkv-311, def:mkv-matrix-of-det-fun, prop:mkv-38-2}
    \lean{MeasureTheory.GibbsMeasure.Markov.matrixOfDetFun_markovDeterminingFun}
    \leanok

    Let $P$ be a positive stochastic matrix and $a \in E$. Then formula (3.7) applied to the
    determining function $g_P$ of (3.11) returns $P$ itself.
\end{theorem}
\begin{proof}
    \leanok

    From (3.11), $Q(x,y) = g_P(a,x,y)/g_P(a,a,y) = P(a,x)P(x,y)/(P(a,a)P(a,y))$. Setting
    $q = 1/P(a,a)$ and $r(x) = P(a,x)$ this reads $Q(x,y)r(y) = q\,r(x)P(x,y)$; summing over $y$
    and using that $P$ is stochastic shows $(q,r)$ is a positive eigenpair of $Q$. By Proposition
    \ref{prop:mkv-38-2} it may be used in (3.7), which then returns $P$.
\end{proof}

\begin{corollary}[Georgii, Theorem (3.5): injectivity of $P \mapsto \gamma_P$]
    \label{cor:mkv-inj}
    \uses{thm:mkv-step2, thm:mkv-311}
    \lean{MeasureTheory.GibbsMeasure.Markov.markovSpecification_injOn}
    \leanok

    If $P, P'$ are positive stochastic matrices with $\gamma_P = \gamma_{P'}$, then $P = P'$.
\end{corollary}
\begin{proof}
    \leanok

    Evaluating the singleton kernels at $i = 0$ on the boundary condition
    $\omega = (x \text{ at } -1, z \text{ elsewhere})$ and using (3.11) gives $g_P = g_{P'}$.
    Now apply Theorem \ref{thm:mkv-step2} to both.
\end{proof}

\begin{theorem}[Georgii, step 4 of the proof of (3.5): equation (3.12)]
    \label{thm:mkv-312}
    \uses{def:mkv-pos-hom-markov, def:markov-spec}
    \lean{MeasureTheory.GibbsMeasure.Markov.eq_312_of_isPositiveHomogeneousMarkovWith, MeasureTheory.GibbsMeasure.Markov.sum_determiningFun_eq_one}
    \leanok

    The determining function $g$ of a positive homogeneous Markov specification $\gamma$ is
    normalised, $\sum_y g(x,y,z) = 1$, and satisfies, for all $x,y,z,a \in E$,
    \[
        \frac{g(x,y,z)}{g(x,a,z)}\cdot\frac{g(a,x,a)}{g(a,a,a)}
          = \frac{g(a,x,y)}{g(a,a,y)}\cdot\frac{g(a,y,z)}{g(a,a,z)}. \qquad (3.12)
    \]
\end{theorem}
\begin{proof}
    \leanok

    Take $\Lambda = \set{1,2}$ and a boundary condition $\omega$ with $\omega_0 = a$,
    $\omega_3 = z$, and write $[xy] = \gamma_\Lambda(\sigma_1 = x, \sigma_2 = y \mid \omega)$.
    Properness of $\gamma$ plus the consistency equations
    $\gamma_\Lambda = \gamma_\Lambda\gamma_{\set{1}} = \gamma_\Lambda\gamma_{\set{2}}$ give
    $[xy] = g(a,x,y)\sum_u [uy] = g(x,y,z)\sum_v [xv]$, whence
    $g(a,x,y)/g(a,a,y) = [xy]/[ay]$ and $g(x,y,z)/g(x,a,z) = [xy]/[xa]$. Both sides of (3.12)
    equal $[xy]/[aa]$. (All the $[xy]$ are strictly positive because $g$ is.)
\end{proof}

\begin{theorem}[Georgii, step 3 of the proof of (3.5): (3.12) $\Rightarrow$ (3.11)]
    \label{thm:mkv-step3}
    \uses{thm:mkv-312, def:mkv-matrix-of-det-fun, thm:mkv-311}
    \lean{MeasureTheory.GibbsMeasure.Markov.markovDeterminingFun_of_eq_312}
    \leanok

    Let $g > 0$ on $E^3$ be normalised, $\sum_y g(x,y,z) = 1$, and satisfy (3.12), and let
    $Q = (g(a,x,y)/g(a,a,y))_{x,y}$. Then for \emph{every} $q > 0$ and every $v \in
    \,]0,\infty[^E$, the matrix $M(x,y) = Q(x,y)v(y)/(q\,v(x))$ has $g$ as its determining
    function (3.11): $g(x,y,z) = M(x,y)M(y,z)/M^2(x,z)$. Taking for $(q,v)$ the Perron pair of
    $Q$ gives the matrix $P$ of (3.7).

    The Lean statement is stronger than Georgii's step 3, which only treats the Perron pair: no
    eigenvector equation is assumed, only positivity of $q$ and $v$.
\end{theorem}
\begin{proof}
    \leanok

    Equation (3.12) rearranges to $g(x,y,z)\,Q(x,a) = g(x,a,z)\,Q(x,y)Q(y,z)$, and the shape of
    $M$ gives $M(x,y)M(y,z)\,q^2v(x) = Q(x,y)Q(y,z)v(z)$, in particular
    $M(x,a)M(a,z)\,q^2v(x) = Q(x,a)v(z)$ since $Q(a,z) = 1$. Combining,
    $g(x,y,z)\,M(x,a)M(a,z) = g(x,a,z)\,M(x,y)M(y,z)$; summing over $y$ and using
    $\sum_y g(x,y,z) = 1$ identifies $M(x,a)M(a,z) = g(x,a,z)M^2(x,z)$, and substituting back
    gives (3.11).
\end{proof}

\begin{theorem}[Georgii, Theorem (3.5): surjectivity of $P \mapsto \gamma_P$]
    \label{thm:mkv-surj}
    \uses{thm:mkv-312, thm:mkv-step3, def:mkv-matrix-of-det-fun, rmk:mkv-comment-32, def:mkv-specification}
    \lean{MeasureTheory.GibbsMeasure.Markov.exists_matrix_eq_markovSpecification, MeasureTheory.GibbsMeasure.Markov.eq_markovSpecification_of_determiningFun}
    \leanok

    Every positive homogeneous Markov specification $\gamma$ on $\mathbb{Z}$ is of the form
    $\gamma = \gamma_P$ for a positive stochastic matrix $P$, namely the matrix computed from the
    determining function of $\gamma$ by (3.7).
\end{theorem}
\begin{proof}
    \leanok

    Let $g$ be the determining function of $\gamma$ and let $P$ be given by (3.7). By Theorem
    \ref{thm:mkv-312} $g$ is normalised and satisfies (3.12), so by Theorem \ref{thm:mkv-step3}
    $g = g_P$ is the determining function of $\gamma_P$ as in (3.11). Thus $\gamma$ and
    $\gamma_P$ have the same singleton kernels, and Remark \ref{rmk:mkv-comment-32} (Theorem
    (1.33)) gives $\gamma = \gamma_P$.
\end{proof}

\begin{theorem}[Georgii, step 1 of the proof of (3.5): existence]
    \label{thm:mkv-step1}
    \uses{def:mkv-chain, def:mkv-specification, thm:mkv-chain-cylinder, thm:mkv-311, def:gibbs-meas}
    \lean{MeasureTheory.GibbsMeasure.Markov.isGibbsMeasure_markovSpecification_stationaryChain}
    \leanok

    For every positive stochastic matrix $P$, the stationary Markov chain $\mu_P$ is a Gibbs
    measure for $\gamma_P$: $\mu_P \in \Gibbs{\gamma_P}$.
\end{theorem}
\begin{proof}
    \leanok

    By Theorem (1.33) for Gibbs measures it suffices that $\mu_P\gamma_{P,\set{i}} = \mu_P$ for
    every site $i$, and by a $\pi$-system argument it suffices to check this on interval
    cylinders $[a,b]$ with $a < i < b$. Theorem \ref{thm:mkv-chain-cylinder} and
    Chapman--Kolmogorov ($\sum_y P(x,y)P(y,z) = P^2(x,z)$) give the three-point identity
    $\mu_P(\sigma_{[a,b]} = \zeta) = g_P(\zeta_{i-1},\zeta_i,\zeta_{i+1})\,
    \mu_P(\sigma_{[a,b]\setminus\set{i}} = \zeta)$, which combined with (3.11) and properness of
    $\gamma_P$ is exactly the required invariance.
\end{proof}

\begin{lemma}[Georgii, step 5 of the proof of (3.5): the Doeblin limit]
    \label{lem:mkv-doeblin-limit}
    \uses{thm:mkv-38-1, lem:mkv-doeblin}
    \lean{MeasureTheory.GibbsMeasure.Markov.tendsto_markovSpecification_Icc_apply_cyl}
    \leanok

    Fix $\Lambda = [l,m]$ with $m = l+d$ and put $\Delta(k) = [l-1-k, m+1+k]$. Then for every
    $\zeta$ and every boundary condition $\omega$,
    \[
        \lim_{k \to \infty} \gamma_{P,\Delta(k)}(\sigma_\Lambda = \zeta_\Lambda \mid \omega)
          = \alpha_P(\zeta_l) P(\zeta_l,\zeta_{l+1}) \cdots P(\zeta_{m-1},\zeta_m).
    \]
\end{lemma}
\begin{proof}
    \leanok

    By Theorem \ref{thm:mkv-38-1} the left-hand side is
    $P^{k+2}(\omega_{l-2-k}, \zeta_l)\,\pi(\zeta)\,P^{k+2}(\zeta_m, \omega_{m+2+k})
    / P^{d+2k+4}(\omega_{l-2-k}, \omega_{m+2+k})$ with
    $\pi(\zeta) = \prod_{j=l}^{m-1}P(\zeta_j,\zeta_{j+1})$. By Lemma \ref{lem:mkv-doeblin} each
    matrix power is within $2\rho^k$ of $\alpha_P$ at the relevant entry. The two boundary sites
    move with $k$, so one estimates the ratio rather than the individual limits: since
    $\alpha_P \ge c > 0$, the numerator and denominator factors at the varying site cancel to
    $1 + O(\rho^k)$, and the remaining factor tends to $\alpha_P(\zeta_l)$.
\end{proof}

\begin{theorem}[Georgii, Theorem (3.5): uniqueness]
    \label{thm:mkv-uniqueness}
    \uses{lem:mkv-doeblin-limit, def:gibbs-meas, def:mkv-specification, lem:gibbs-meas-tfae}
    \lean{MeasureTheory.GibbsMeasure.Markov.isGibbsMeasure_apply_cyl, MeasureTheory.GibbsMeasure.Markov.eq_of_isGibbsMeasure}
    \leanok

    Let $P$ be a positive stochastic matrix. Every $\mu \in \Gibbs{\gamma_P}$ satisfies
    $\mu(\sigma_\Lambda = \zeta_\Lambda) = \alpha_P(\zeta_l)\prod_{j=l}^{m-1}
    P(\zeta_j,\zeta_{j+1})$ for every interval $\Lambda = [l,m]$; consequently
    $\Gibbs{\gamma_P}$ has at most one element.
\end{theorem}
\begin{proof}
    \leanok

    For $\mu \in \Gibbs{\gamma_P}$ one has $\mu\gamma_{P,\Delta(k)} = \mu$, so
    $\mu(\sigma_\Lambda = \zeta_\Lambda) = \int \gamma_{P,\Delta(k)}(\sigma_\Lambda =
    \zeta_\Lambda\mid\omega)\,\mu(d\omega)$ for all $k$. Dominated convergence and Lemma
    \ref{lem:mkv-doeblin-limit} identify the limit as the chain value. Interval cylinders
    determine a cylinder of an arbitrary finite volume by summing over the extra sites, and the
    cylinders form a $\pi$-system generating $\mathcal{F}$, so $\mu$ is determined.
\end{proof}

\begin{theorem}[Georgii (3.6)]
    \label{thm:mkv-36}
    \uses{thm:mkv-step1, def:mkv-chain, def:mkv-specification, def:cond-exp-ker}
    \lean{MeasureTheory.GibbsMeasure.Markov.markovSpecification_apply_cyl_eq_cond,
      MeasureTheory.GibbsMeasure.Markov.isCondExp_markovSpecification_stationaryChain}
    \leanok

    For every $\Lambda \in \Fin$, every $\zeta$ and \emph{every} --- not merely almost every ---
    boundary condition $\omega$,
    \[\gamma_{P,\Lambda}(\sigma_\Lambda = \zeta \mid \omega) =
      \mu_P(\sigma_\Lambda = \zeta \mid \sigma_{\partial\Lambda} = \omega_{\partial\Lambda}),\]
    conditioning on the finite boundary $\partial\Lambda$ of (3.4) alone: Georgii's statement
    verbatim, as \texttt{markovSpecification\_apply\_cyl\_eq\_cond}, with the conditioning event
    of positive $\mu_P$-measure (\texttt{stationaryChain\_cyl\_pos}), so the right-hand side is
    the elementary conditional probability.  The kernel-level companion --- $\gamma_{P,\Lambda}$
    is a version of the conditional distribution of $\mu_P$ given the whole exterior
    $\sigma$-algebra $\mathcal{F}_{\mathbb{Z}\setminus\Lambda}$ --- is
    \texttt{isCondExp\_markovSpecification\_stationaryChain}.  For an interval
    $\Lambda = [a,b]$ the boundary is $\set{a-1, b+1}$ (\texttt{boundary\_Icc}) and the
    right-hand side is evaluated by Comment (3.8)(1) as Theorem \ref{thm:mkv-38-1-single}.
\end{theorem}
\begin{proof}
    \leanok

    Both sides are ratios of Boltzmann weights of the Markov potential: the left-hand side by
    the definition of the Markov specification, the right-hand side by the single-site exchange
    identity for $\mu_P$
    (\texttt{stationaryChain\_cyl\_update\_mul\_boltzmannFactor}), inducting over the sites of
    $\Lambda$.  The kernel-level version is Theorem \ref{thm:mkv-step1}.
\end{proof}

\begin{theorem}[Georgii (3.5)]
    \label{thm:mkv-35}
    \uses{thm:mkv-step1, thm:mkv-uniqueness, thm:mkv-surj, cor:mkv-inj, def:mkv-chain, def:mkv-pos-hom-markov, def:gibbs-meas}
    \lean{MeasureTheory.GibbsMeasure.Markov.georgii_3_5, MeasureTheory.GibbsMeasure.Markov.gibbsMeasure_eq_singleton, MeasureTheory.GibbsMeasure.Markov.existsUnique_isGibbsMeasure}
    \leanok

    The relation $\Gibbs{\gamma} = \set{\mu_P}$ establishes a one-to-one correspondence
    $\gamma \leftrightarrow P$ between the set of all positive homogeneous Markov specifications
    on $\mathbb{Z}$ and the set of all stochastic matrices on $E$ with non-vanishing entries. For
    a given $P$ the corresponding $\gamma$ is determined by (3.6), and conversely $P$ is
    expressed in terms of the determining function $g$ of $\gamma$ by (3.7).
\end{theorem}
\begin{proof}
    \leanok

    Georgii's five steps are, in Lean: step 1, $\mu_P \in \Gibbs{\gamma_P}$ and $\gamma_P$ is a
    specification (Theorem \ref{thm:mkv-step1}; properness and consistency are free here, since
    $\gamma_P$ is Gibbsian for $\Phi^P$ by construction), with the determining function (3.11)
    (Theorem \ref{thm:mkv-311}); step 2, $P$ is recovered from $g_P$ by (3.7), whence injectivity
    (Theorem \ref{thm:mkv-step2}, Corollary \ref{cor:mkv-inj}); steps 3 and 4, every determining
    function satisfies (3.12) and (3.12) forces the shape (3.11), whence surjectivity (Theorems
    \ref{thm:mkv-312}, \ref{thm:mkv-step3}, \ref{thm:mkv-surj}); step 5, the ergodic theorem for
    Markov chains gives $\Gibbs{\gamma_P} \subseteq \set{\mu_P}$ (Theorem
    \ref{thm:mkv-uniqueness}).
\end{proof}

\begin{corollary}[Georgii (3.9)]
    \label{cor:mkv-39}
    \uses{thm:mkv-surj, def:mkv-potential, def:mkv-pos-hom-markov}

    A specification $\gamma$ is a positive homogeneous Markov specification if and only if
    $\gamma$ is Gibbsian for some homogeneous nearest-neighbour potential $\Phi$, i.e.\ one with
    $\Phi_{\set{i}} = \varphi_1(\sigma_i)$ and $\Phi_{\set{i,i+1}} =
    \varphi_2(\sigma_i,\sigma_{i+1})$.

    Only the forward implication is formalized, and it is exactly Theorem \ref{thm:mkv-surj}:
    $\gamma = \gamma_P$ is Gibbsian for the homogeneous nearest-neighbour potential
    $\Phi^P = -\log P$ of Definition \ref{def:mkv-potential}. The converse --- that the Gibbsian
    specification of an \emph{arbitrary} homogeneous nearest-neighbour pair
    $(\varphi_1,\varphi_2)$ is a positive homogeneous Markov specification --- is not in the
    library, and neither is the consequence Georgii draws from it with (2.35)(b): that for a
    fixed $\alpha \in \Prob(E,\mathcal{E})$ the relation $\Gibbs{\Phi} = \set{\mu_P}$ is a
    bijection between the $\alpha$-normalised homogeneous nearest-neighbour potentials and the
    positive stochastic matrices.
\end{corollary}

\begin{remark}[The half-line variant]
    \label{rmk:mkv-half-line}
    \uses{thm:mkv-35}

    Georgii notes that Theorem (3.5) has a straightforward modification for $S = \mathbb{N}$, with
    $\mu_P$ replaced by the Markov chain with transition matrix $P$ started at time $0$ in a fixed
    state $a \in E$. This variant is not formalized; the whole file is written for
    $S = \mathbb{Z}$.
\end{remark}
